\chapter{Налагодження машини}
\chaptermark{Налагодження машини}
\label{sec:debugging}

\section{Як налагоджують машину без схеми}

Інженер, якому дали робочу плату без схеми, налагоджує її чотирма прийомами. Він ставить \emph{щуп} і дивиться, що йде проводом. Він \emph{подає тестовий сигнал} і дивиться, що змінилося. Він \emph{вимикає} шматок схеми і дивиться, що перестало працювати. І він \emph{читає регістри керування}, щоб дізнатися, в якому режимі працює плата. Уся ця книга --- спроба прочитати схему мозку, і в цьому розділі подивимося, якими з чотирьох прийомів інженери вже вміють користуватися і що кажуть їм попередні розділи про те, чого чекати.

Найпомітніший сьогодні приклад такого налагодження --- інтерфейси мозок--комп'ютер: пристрої, які читають імпульси нейронів і перетворюють їх на команди для зовнішніх машин або, навпаки, подають у мозок сигнали ззовні. Їм і присвячено більшу частину розділу, і насамперед тому, що робить компанія Neuralink.

\section{Щуп: що чує електрод}

Щуп мозку --- тонкий електрод, уведений у тканину. Він чує імпульси нейронів, що лежать у радіусі порядку сотні мікрометрів, --- зазвичай від одного до кількох. Імпульс на електроді --- сплеск у десятки--сотні мікровольтів тривалістю близько мілісекунди, і за його формою можна відрізнити сусідні нейрони один від одного. Найкраще чутно великі \nt{GLUT}-нейрони: їх у корі більшість (таблиця~\ref{tab:types}), і їхні струми сильніші.

Одразу видно головну трудність. Добрий щуп сьогодні --- це сотні або тисячі електродів, а нейронів у мозку $8.6\cdot10^{10}$. Ми підслуховуємо приблизно одну стомільйонну машини. Чому ж із такої мізерної вибірки взагалі можна щось зрозуміти? Відповідь дав розділ~\ref{sec:code}. Сусідні нейрони шумлять узгоджено, і усереднення за групою впирається в межу (твердження~\ref{prop:pooling}): сотня нейронів несе майже стільки ж, скільки тисяча. До того ж задача, яку розв'язує рухова кора, маловимірна: у руки небагато суглобів (розділ~\ref{sec:muscles}), і активність сотень нейронів, що нею керують, укладається в простір лише з одного-двох десятків спільних змінних~\cite{Gallego2017}. Щуп на сотню нейронів чує майже всі ці змінні. Якби мозок зберігав незалежне повідомлення в кожному нейроні, таке підслуховування було б марним.

\section{Потік даних: стиснення просто на щупі}

Щуп видає величезний потік. Щоб побачити форму імпульсу, кожен електрод треба оцифровувати приблизно $20$ тисяч разів на секунду з точністю близько десяти біт. Тисяча електродів дає
\begin{empheq}[box=\keybox]{equation}
\begin{aligned}
R_{\text{сирий}} \;&=\; N\,f_s\,b \;\approx\; 10^3\cdot2\cdot10^4\cdot10 \;\approx\; 2\cdot10^8\ \text{біт/с},\\
R_{\text{імп}} \;&\approx\; N\,r\,\log_2\frac{e}{r\,\Delta t} \;\approx\; 8\cdot10^4\ \text{біт/с}.
\end{aligned}
\label{eq:probedata}
\end{empheq}
Друга оцінка --- пропускна здатність потоку імпульсів~\eqref{eq:spikeentropy} за $r=10$ імпульсів на секунду і точності $\Delta t=1\ms$: усе, що в ньому насправді є, займає у дві з половиною тисячі разів менше. Для вживленого пристрою це вирішальна різниця: сирий потік не можна передати радіо крізь шкіру з тим живленням, яке можна дозволити собі всередині голови, не нагріваючи тканини. Тому вживлений пристрій має виділяти імпульси просто на кристалі поруч з електродами й передавати назовні лише події --- номер каналу й момент імпульсу. У першому описі системи Neuralink до цього ще не дійшли: кристали підсилюють і оцифровують сигнали, увесь сирий потік іде кабелем, а імпульси виділяє програма ззовні; стиснення на кристалі й бездротовий зв'язок названо там планом~\cite{Musk2019}.

Це знайомий прийом. Око стискає сигнал у сто разів, перш ніж надіслати його довгим кабелем (розділ~\ref{sec:sensors}); подієва камера передає не кадри, а зміни. Щуп, щоб умістися в провід, змушений робити те саме, що робить сам мозок: говорити імпульсами і лише про новини.

\section{Що робить Neuralink}

Пристрій Neuralink --- щуп, зібраний під ці обмеження~\cite{Musk2019}. Замість жорсткої ґратки голок --- багато гнучких ниток, тонших за волосину, з електродами вздовж кожної: у першому описі системи --- до $3072$ електродів на $96$ нитках, по $32$ на нитку, а в пристрої, вживленому людям, за повідомленнями компанії, близько тисячі каналів на $64$ нитках. Гнучкі нитки менше ушкоджують тканину й краще переносять те, що мозок усередині черепа весь час трохи зсувається. Нитки вводить робот, оминаючи судини. У першому описі, як сказано вище, сирий потік~\eqref{eq:probedata} ішов назовні кабелем. Пристрій для людей, за повідомленнями компанії, влаштований інакше: корпус повністю закритий шкірою, імпульси виділяються всередині, дані йдуть радіо, живлення надходить бездротово.

Завдання першого пристрою --- читання. Нитки ставлять у ту частину кори, яка планує рухи руки, і декодер перетворює імпульси на рух курсора на екрані. Людина з паралічем рук керує комп'ютером, уявляючи рух. Сама ідея не нова: інтерфейси на жорстких ґратках із сотні електродів давно дозволяли керувати курсором~\cite{Hochberg2006}, писати текст, уявляючи рухи руки під час письма~\cite{Willett2021}, і навіть перетворювати спроби говорити на текст на екрані зі швидкістю близько шістдесяти слів на хвилину~\cite{Willett2023}. Нове в Neuralink --- інженерія: кількість каналів, бездротовість, закритий корпус і роботизоване введення, тобто спроба зробити щуп, який можна носити роками поза лабораторією. Наскільки нам відомо, результати випробувань на людях публікувала переважно сама компанія, а не рецензовані статті; компанія повідомляла, зокрема, що частина ниток після введення змістилася і кількість робочих каналів зменшилася. Для налагодження це звичайна неприємність: щуп, який рухається відносно плати, чує вже інші проводи.

Другий напрям компанії --- запис: пристрій для зору, який подає сигнали в зорову кору людини, що втратила зір. Про нього нижче, у підрозділі про запис.

\section{Читання: декодер як отримувач}

Декодер інтерфейсу --- це отримувач у сенсі твердження~\ref{prop:reader}: він сам вирішує, яку частину повідомлення прочитати. Практично всі інтерфейси читають \emph{частоти груп}: рахують імпульси кожного каналу у вікнах по кілька десятків мілісекунд і додають їх із вагами, дібраними так, щоб виходив задуманий рух. Це частотний код групи з підрозділу~\ref{sec:rate}, і його обрано не випадково: коли на вікно припадає кілька імпульсів, частота одного нейрона не визначена, а сума за сотнею вже працює.

Скільки інформації вдається витягти? Письмо «уявною рукою» йшло зі швидкістю близько дев'яноста знаків на хвилину~\cite{Willett2021}, тобто півтора знака на секунду: навіть за п'яти бітів на знак це менше ніж вісім біт на секунду. Керування курсором, за повідомленнями Neuralink, дає порядку десяти біт на секунду. А верхня межа~\eqref{eq:spikeentropy} для \emph{одного} нейрона --- десятки біт на секунду, для сотні --- тисячі. Інтерфейс витягає з підслуханих нейронів малу частку того, що вони могли б нести. Вузьке місце --- не провід, а те, скільки незалежних змінних несе група (твердження~\ref{prop:pooling}) і наскільки добре декодер розуміє код, який, як ми бачили в розділі~\ref{sec:code}, до кінця не розшифровано.

Є й друга особливість, знайома з розділу~\ref{sec:motorlearning}. Декодер і мозок учаться один в одного. Мозок бачить, куди поїхав курсор, отримує помилку і підлаштовує свої команди --- те саме навчання рухів проводом помилки. А інженери час від часу заново добирають ваги декодера, бо нейрони під нитками змінюються і зв'язки в мережі перенавчаються. Виходять два учні, які підлаштовуються один під одного; щоб така пара не втратила стійкості, швидкості їхнього навчання зручно розвести: нехай один учиться швидко, а другий повільно.

\begin{keyblock}
\begin{proposition}[Чому щуп на тисячу нейронів працює]
\label{prop:probe}
Інтерфейс, що підслуховує мізерну частку нейронів, може керувати рухом тому, що активність групи маловимірна і її шуми узгоджені: кілька сотень нейронів несуть майже всі спільні змінні. З тієї самої причини інтерфейс витягає лише кілька біт на секунду --- стільки, скільки незалежних змінних у задачі, а не стільки, скільки вміщує потік імпульсів. Роботу інтерфейсів виміряно; наскільки кращим стане декодер, коли код буде зрозуміло, --- відкрите питання.
\end{proposition}
\end{keyblock}

\section{Запис: важче, ніж читання}

Подати сигнал у машину важче, ніж прочитати. Електрод, через який пропускають струм, збуджує не один нейрон, а всі нейрони й проводи навколо себе, без розбору типу й знака. Записати ним код, у якому важливо, \emph{який} нейрон і \emph{коли} спрацював (розділ~\ref{sec:code}), не можна: можна лише грубо задати, \emph{де} і \emph{наскільки сильно}.

Для зору цього частково вистачає. Струм через електрод у зоровій корі людина бачить як світну точку в певному місці поля зору; коли точки запалювали одночасно, до шістнадцяти разом, людина, що втратила зір, упізнавала деякі літери й розрізняла межі предметів~\cite{Fernandez2021}. Це код місцем, і його записати можна. Але згадайте розділ~\ref{sec:sensors}: око надсилає в мозок не пікселі, а стиснений опис --- краї, зміни, посвітлішання й потемніння, окремими проводами. Запис «точок світла» в кору --- це запис пікселів у машину, яка чекає на вході зовсім іншого коду. Тому штучний зір поки грубіший, ніж можна було б чекати з кількості електродів. Є й випробувальні сигнали, яким електрод не потрібен: зорові ілюзії подають через око незвичний вхід і виводять машину зі звичайного режиму, а кожна помилка вказує на свій механізм (підрозділ~\ref{sec:illusions}).

\section{Чого щуп не бачить}

Електрод чує адресну шину --- імпульси \nt{GLUT}-нейронів і, гірше, дрібних \nt{GABA}. Але в розділах~\ref{sec:serocomputer}--\ref{sec:acet} ми бачили, що режим роботи всієї машини задають широкомовні регістри: концентрації \nt{SERO}, \nt{DOPA}, \nt{NORA}, \nt{ACET}. Їх електрод не чує: нейронів-джерел мізерно мало, і їхній сигнал --- не імпульс поруч зі щупом, а концентрація в об'ємі. Налагоджувач, який бачить шину даних, але не бачить регістрів керування, завжди дивуватиметься: той самий вхід даватиме різні відповіді, бо десь змінилося $g$, $a$ чи $\tau_\gamma$. Концентрації можна вимірювати хімічними датчиками просто в тканині~\cite{Bunin1998}, але в інтерфейсах такі датчики поки не використовують. Зовсім поза полем зору щупа лишається й другий, повільний вихід машини --- гормони в крові (розділ~\ref{sec:chemoutput}): їх вимірюють у пробах крові, а не електродом.

Не бачить щуп і ваг --- а саме в них зберігається майже вся пам'ять і все вивчене. Їх можна лише вгадувати з того, як змінюються відповіді. І не бачить він болю так, як його відчуває людина: у розділі~\ref{sec:pain} ми бачили, що силу сигналу тривоги задає сама машина --- воротами в спинному мозку й регуляторами згори, тож із датчика не можна прочитати, наскільки боляче.

\section{Підсумок: налагодження і відкриті питання}

\begin{table}[t]
\centering
\footnotesize
\setlength{\tabcolsep}{4pt}
\renewcommand{\arraystretch}{1.15}
\begin{tabular}{>{\raggedright}p{2.2cm}>{\raggedright}p{3.6cm}>{\raggedright}p{3.4cm}>{\raggedright\arraybackslash}p{4.0cm}}
\toprule
Прийом налагодження & Що робить інтерфейс & Що пояснює книга & Що відомо \\
\midrule
щуп & чує сотні--тисячі нейронів & маловимірність і узгоджений шум (розд.~\ref{sec:code}) & виміряно \\
стиснення на щупі & передає події замість сирого сигналу & пропускна здатність потоку імпульсів~\eqref{eq:spikeentropy} & інженерно розв'язано \\
читання & частоти груп $\to$ рух, текст, мовлення & декодер --- отримувач, частотний код групи & працює; кілька біт на секунду \\
спільне навчання & мозок і декодер підлаштовуються один під одного & провід помилки (розд.~\ref{sec:motorlearning}) & спостерігається; правила підлаштування --- предмет досліджень \\
запис & струм запалює точки в полі зору & код місцем записати можна, часовий --- ні (розд.~\ref{sec:sensors}) & точки виміряно; корисного зору --- поки немає \\
читання регістрів & майже не робиться & широкомовні канали (розд.~\ref{sec:serocomputer}--\ref{sec:acet}) & хімічні датчики є в дослідах \\
\bottomrule
\end{tabular}
\caption{Налагодження машини: що вміють інтерфейси мозок--комп'ютер і що про них кажуть розділи книги.}
\label{tab:debugging}
\end{table}

Таблиця~\ref{tab:debugging} підсумовує розділ. Інтерфейси мозок--комп'ютер уже вміють ставити щуп на адресну шину й читати з неї кілька біт на секунду, і те, що це взагалі працює, пояснюється маловимірністю й узгодженим шумом із розділу~\ref{sec:code}. Записувати в машину вони вміють лише грубо, бо її вхідний код стиснений і адресований окремим проводам. А регістри керування й ваги --- те, що робить машину здатною навчатися й налаштовуватися, --- щупові поки майже не видно.

Кожне відкрите питання розділу~\ref{sec:synthesis} --- це й задача для налагоджувача. Який код читати, вирішиться, коли щупи стануть щільнішими, а декодери навчаться користуватися часом імпульсів, а не лише частотами. Як навчається багатошарова мережа, можна перевірити, поставивши щупи одразу в кілька шарів і дивлячись, як змінюються їхні відповіді під час навчання. Що передають широкомовні канали, стане видно, коли до електричних щупів додадуться хімічні. Ця книга --- спроба прочитати схему машини за її поведінкою і за законами, які знає будь-який інженер. Налагоджувачі, що будуються зараз, дадуть змогу перевірити цю схему щупом.

\section{Задачі}

\begin{exercise}[\lvlA]
\label{ex:17:1}
\emph{Бюджет радіоканалу.} Передавання крізь шкіру коштує $1$~нДж на біт, передавачеві можна витрачати не більше $10$~мВт. Яка потужність потрібна для сирого потоку~\eqref{eq:probedata} за $N=1000$ електродів і для потоку імпульсів? Якої енергії на біт вимагав би сирий потік?
\end{exercise}

\begin{exercise}[\lvlA]
\label{ex:17:2}
\emph{Скільки біт у сотні нейронів.} Декодер додає імпульси $N$ нейронів за вікна $50\ms$ за $r=20$~імп/с, попарна кореляція шуму $c=0.1$, сигнал змінює частоту групи на $30\%$ (середньоквадратично). Оцініть швидкість $\tfrac12\log_2(1+\text{SNR})$ за вікно за $N=10$, $100$, $1000$ і $N\to\infty$. А за $c=0$, $N=1000$?
\end{exercise}

\begin{exercise}[\lvlB]
\label{ex:17:3}
\emph{Швидкість інтерфейсу-клавіатури.} Інтерфейс півтора раза на секунду обирає один із $N=32$ знаків: задуманий з імовірністю $P$, помилки рівноймовірні. Скільки біт на секунду він передає за $P=0.99$, $0.94$, $0.8$? Скільки знаків на секунду треба за $P=0.94$ для $10$~біт/с?
\end{exercise}

\begin{exercise}[\lvlB]
\label{ex:17:4}
\emph{Моделювання: декодер групи.} $N$ нейронів із $\lambda_i=[b+m\,\mathbf v\cdot\mathbf p_i]_+$, $b=20$, $m=15$~імп/с, вікна $50\ms$, $\mathbf v$ зі стандартним відхиленням $0.5$ за компонентою; усі бачать спільний шум, $\mathbf v+\boldsymbol\xi$, $\sigma_\xi=0.2$. Змоделюйте незміщений декодер вектором групи. За якого $N$ обидва шуми рівні і скільки біт на компоненту за вікно дає $N\to\infty$?
\end{exercise}

\begin{exercise}[\lvlB]
\label{ex:17:5}
\emph{Два учні.} Мозок видає команду $m=b\,v^*$, декодер --- $v=d\,m$, $\langle v^{*2}\rangle=1$. Обидва після кожної спроби роблять крок градієнта за $\tfrac12(v-v^*)^2$ зі швидкістю $\eta$. Змоделюйте з $b=1.2$, $d=1$ за $\eta=0.8$, $1.1$, $1.5$, $2$. Кожен поодинці стійкий за $\eta<2$; за якого $\eta$ стійка пара?
\end{exercise}

\begin{exercise}[\lvlB]
\label{ex:17:6}
\emph{Проєктування: конвеєр щупа.} $1024$ канали по $20$~тис. відліків на секунду, нейрони $10$~імп/с, радіо $1$~нДж/біт. Який поріг за білим шумом відліків дає не більше $0.1$~Гц хибних імпульсів на канал? Передаємо числа імпульсів за $20\ms$ по $2$~біти: яка потужність і у скільки разів вона менша, ніж для сирих $10$-бітних відліків?
\end{exercise}

\begin{exercise}[\lvlC]
\label{ex:17:7}
\emph{Межа швидкості інтерфейсу.} Оцініть $R\le DW\log_2(1+\text{SNR})$ для $D=10$ змінних зі смугою $W=5$~Гц за $\text{SNR}\approx0.9$ (шум, схожий на сигнал) і $90$ (незалежний шум $1000$ нейронів). Виміряно близько $10$~біт/с. Який дослід відрізнить два пояснення розриву: шум, схожий на сигнал, і незнання коду?
\end{exercise}
